\begin{figure*}[!htb]
\small
\begin{tcolorbox}[
    colback=white, 
    colframe=blue3, 
    boxrule=0.5mm,
    arc=0mm,
    outer arc=0mm,
    title=Limited Functionality Tools: Iterative - Generation Prompt
]
Task Description:\\
You need to go ahead and generate samples of tool usage with as much diversity as possible, based on the current dimension, the requirements requirements of the task, and the sample references provided.\\

Current dimension:\\
Iterative Dimension - The task requires repeated invocations of the same tool, with multiple iterations to get the final result\\

Task Requirement:\\
- Generate new samples taking into account as much as possible the diversity of user query and available tools.
- First ensure that the task description is a goal that requires the use of multiple tools to accomplish, and provide the corresponding tools in the list of available tools using numerical numbers. Then make sure that one tool is called more than twice in a row during the planning phase, and that this row of calls is reflected in the task planning.\\
- Use "// Raw Input //" and "// Raw output //" to label the inputs and outputs of the samples.\\
- The samples need to consider a realistic possible scenario as much as possible. Multiple tools need to be available in the list of available tools in the raw input. Note that in the current dimension, the raw output has multiple iterative calls to the same tool in the planning phase, with at least two or more steps in between, corresponding to the iterative call dimension.\\

At build time, you need to:\\
First give a list of overall tasks and then express specific task requirements separately. In the list of available tools, make sure to provide tools that can fulfill these requirements, while including a small number of tools that do not directly help the task to ensure task complexity and tool diversity. Second, by placing certain limitations on the tool capabilities that result in multiple invocations to accomplish a given requirement, and making these limitations and requirements explicit in the task description and the list of provided tools. Finally, the corresponding task planning process is completed in the output based on these requirements.\\

Sample Reference:\\
==== Reference Sample Start ====\\
// Raw Input //\\
\colorbox{blue3!50}{\texttt{[sample\_input]}}\\

// Raw output //\\
\colorbox{blue3!50}{\texttt{[sample\_outputl]}}\\
==== Reference Sample End ====\\

Now, please combine the reference sample and the corresponding build steps, try to select tasks from different domains and scenarios, and use different types of tools to make the new sample different from the reference sample in terms of task content and context. We start generating new single samples corresponding to the dimensions.\\
Iterative dimension generation sample (one):\\
==== Sample Start ====
\end{tcolorbox}
\caption{Generation prompt for limited functionality tools scenario, iterative subtask.}
\label{fig:iter-prompt}
\end{figure*}